% Chapter 3, section 8: Last value
% Write in plain LaTeX. Code of a guide step: \guide{<step>}; commands: \begin{shell}...\end{shell};
% a screenshot: \screenshot{file}{caption}; a gap in the draft: \placeholder{...}.
% See book/latex/queen.sty for the full list.
%
% Draft by Claude, 2026-10-24. The code is cut from examples/tutorials/js/book-use-cases.mjs, which
% checks every result quoted here against a 2.7.0 node.

\section{Last value}

The availability of a hotel changes many times a minute. A screen that shows it needs only the
current value, and handling every old update on the way there is wasted work.

A consumer group with \texttt{.conflation()} gets only the newest message of each partition, and its
cursor moves past the older ones.

\queenfig{last-value}{Conflation: the newest value of each partition, for this group only.}

\guide{uc-latest}

Three updates were pushed (12, 11 and 9 rooms). The screen got 9, and the older two were never
delivered to it.

\why{for a value that changes, only the last one matters. Conflation lets a slow consumer catch up
in one step instead of replaying every update. Other groups of the same queue still get every
message.}

One thing to know: use it only where an old message has no meaning of its own. For events that
each count, such as payments, every message must be handled.
